% uid: ""
% school: "lycee"
% level: "1ere"
% chapter: "Activités algébriques"
% title: "Trois nombres de somme nulle"
% difficulty: "Difficile"
% tags: ["identites-remarquables", "factorisation", "cubes", "equations"]
% provenance: "own"
% batch: "2026-11-22-revision-1ere-A"

% ---STATEMENT---
\begin{enumerate}
  \item Développer $(a+b+c)(a^2+b^2+c^2-ab-bc-ca)$. En déduire que, si $a+b+c=0$, alors $a^3+b^3+c^3=3abc$.
  \item Calculer, sans calculatrice, $17^3-9^3-8^3$ puis $20^3-13^3-7^3$.
  \item Factoriser $(x-y)^3+(y-z)^3+(z-x)^3$. En déduire la valeur de $\dfrac{(a-b)^3+(b-c)^3+(c-a)^3}{(a-b)(b-c)(c-a)}$ pour $a$, $b$, $c$ deux à deux distincts.
  \item Résoudre dans $\mathbb R$ l'équation $(x-1)^3+(2x-3)^3+(4-3x)^3=0$.
\end{enumerate}

% ---SOLUTION---
\begin{prop}[Développement]
Pour développer un produit de deux sommes, on multiplie \textbf{chaque} terme de la première par \textbf{chaque} terme de la seconde. Un produit est nul si et seulement si \ib{l'un de ses facteurs est nul}.
\end{prop}
\begin{enumerate}
  \item On multiplie la deuxième parenthèse par $a$, puis par $b$, puis par $c$ :
  \[a\,(\cdots)=a^3+ab^2+ac^2-a^2b-abc-a^2c\]
  \[b\,(\cdots)=a^2b+b^3+bc^2-ab^2-b^2c-abc\]
  \[c\,(\cdots)=a^2c+b^2c+c^3-abc-bc^2-ac^2\]
  On additionne : les termes $ab^2$, $ac^2$, $a^2b$, $a^2c$, $bc^2$, $b^2c$ se détruisent deux à deux, et il reste $-3abc$ :
  \[\ib{$(a+b+c)(a^2+b^2+c^2-ab-bc-ca)=a^3+b^3+c^3-3abc$}\]
  Si $a+b+c=0$, le membre de gauche est nul, donc :
  \[a^3+b^3+c^3-3abc=0\implies\ib{$a^3+b^3+c^3=3abc$}\]

  \item $17+(-9)+(-8)=0$ : on applique la question 1 avec $a=17$, $b=-9$, $c=-8$.
  \[17^3-9^3-8^3=3\times17\times(-9)\times(-8)=3\times17\times72=\ib{$3672$}\]
  $20+(-13)+(-7)=0$ :
  \[20^3-13^3-7^3=3\times20\times(-13)\times(-7)=3\times20\times91=\ib{$5460$}\]

  \item Les trois nombres $x-y$, $y-z$ et $z-x$ ont pour somme
  \[(x-y)+(y-z)+(z-x)=0\]
  Donc, d'après la question 1 :
  \[(x-y)^3+(y-z)^3+(z-x)^3=\ib{$3(x-y)(y-z)(z-x)$}\]
  Avec $x=a$, $y=b$, $z=c$ (le dénominateur n'est pas nul car $a,b,c$ sont distincts) :
  \[\dfrac{(a-b)^3+(b-c)^3+(c-a)^3}{(a-b)(b-c)(c-a)}=\dfrac{3(a-b)(b-c)(c-a)}{(a-b)(b-c)(c-a)}=\ib{$3$}\]

  \item On additionne les trois nombres qui sont élevés au cube :
  \[(x-1)+(2x-3)+(4-3x)=0\]
  La question 1 donne :
  \[(x-1)^3+(2x-3)^3+(4-3x)^3=3\,(x-1)(2x-3)(4-3x)\]
  L'équation devient $3\,(x-1)(2x-3)(4-3x)=0$, soit :
  \[x=1\qquad\text{ou}\qquad x=\dfrac32\qquad\text{ou}\qquad x=\dfrac43\]
  \[S=\ib{$\left\{1\,;\dfrac43\,;\dfrac32\right\}$}\]
\end{enumerate}
\textbf{Idée :} repérer que les trois bases ont une somme nulle, avant de développer quoi que ce soit.

\textbf{Piège :} appliquer $a^3+b^3+c^3=3abc$ sans vérifier que $a+b+c=0$.
